
\section{Methods}
\label{sec:methods}

All computation is implemented in Python~3.9+ as the open-source package
\code{gcrisk}. The pipeline is a
sequence of composable modules: spectrum generation, slab transport, organ
dose, cancer risk (REID), SEP acute hazard, and uncertainty propagation, plus
a separate proton relative biological effectiveness (RBE) self-consistency
check. Full equations, parameter
tables, and derivations for every subsection below are given in
Supplementary Methods (Sections~S1.1--S1.9); only the essential physics and
the numbers needed to interpret the Results are summarized here.

\textbf{Trajectory and solar modulation.}
\label{sec:phi}
A Hohmann-transfer trajectory
(Earth 1.0~AU to Mars 1.524~AU, 259-day transit) sets the mission timeline.
Solar modulation potential $\phi(t)$ comes from the Usoskin et~al.\ monthly
reconstruction (real neutron-monitor data back to 1936); for the MSL transit
epoch, $\phi \approx 481\pm50$~MV. Neutron-monitor inter-station offsets
($\pm50$--$100$~MV) are propagated as an LHS parameter (\S\ref{sec:uq}).

\textbf{GCR spectrum.}
\label{sec:calib}
Solar modulation of GCR flux uses the Gleeson--Axford force-field
approximation applied to per-species local interstellar spectra (LIS):
protons and helium come from Vos \& Potgieter \cite{vos2015}, and C/O/Si/Fe from
the HelMod model of Boschini et al.\ \cite{boschini2020}. Inter-species
flux ratios are constrained to ACE/CRIS \citep{george2009} and PAMELA
\citep{adriani2014} measurements at the MSL epoch. A single global scale
is then fit to match the MSL/RAD absorbed dose of $1.84\mGyday$ at
$16\gcm$ aluminum \citep{zeitlin2013}. This is a one-free-parameter calibration
whose species mix is externally anchored, replacing an earlier
six-free-parameter fit with no compositional constraint. The resulting
HZE fraction ($13.0\%$ at $16\gcm$ Al) is $\sim\!1.7\times$ the NSRL
reference field \citep{slaba2016}, versus $\sim\!4\times$ under the prior
calibration; this residual excess is attributed to the CSDA transport
model not capturing nuclear fragmentation (Supplementary \S{}S1.3.3).

\textbf{Transport.} Ionization energy loss uses CSDA with Barkas
effective-charge scaling on NIST
PSTAR/ASTAR proton stopping powers. Nuclear interactions are modeled as
exponential fluence attenuation with a Bradt-Peters geometric mean free path
\citep{bradtpeters1950}, the same form underlying HZETRN's tabulated mean
free paths \citep{wilson1991}. This does not track secondary fragment
production, which is the model's primary physical limitation, most relevant
above $\sim\!30\gcm$ for Fe/Si. Secondary neutron dose equivalent uses a
pre-computed lookup calibrated to the HZETRN benchmark of Slaba et al.\
\cite{slaba2014} ($\pm25$--$30\%$, via LHS).

\textbf{Organ dose and quality factor.} Absorbed dose and ICRP-60 dose
equivalent \citep{icrp60} are computed at 11 organ-specific tissue depths
from ICRP Publication~89 \citep{icrp89}, weighted by ICRP-60 tissue factors.
The field effective quality factor is $\Qeff = \sum H_{k,s}/\sum D_{k,s}$
(numeric values and the ion-vs.-neutron breakdown: \S\ref{sec:let}).

\textbf{Cancer risk (REID).} REID is computed from the Cucinotta et al.\
\cite{cucinotta2013} organ-specific ERR/EAR model (NASA/TP-2013-217375),
using interpolated US sex-specific survival and baseline-mortality lookup
tables as a practical approximation to the full NASA/Cucinotta framework.
Point estimates use a fixed dose and dose-rate effectiveness factor
(DDREF) of $1.75$; the uncertainty ensemble
instead samples DDREF~$\sim$~Uniform$[1,2]$. REID's distribution is
right-skewed because \code{Q\_factor} (\S\ref{sec:uq}), the dominant tail
driver, is drawn from a heavy-tailed lognormal; the effective quality
factor is clipped at the ICRP-60 physical maximum of $30$, and the
resulting REID sample is separately clipped to $[0,1]$ since it represents
a probability (Supplementary \S{}S1.6).

\textbf{SEP acute risk.} Event fluence spectra use the Band function
\citep{tylka2006} for three canonical events (August~1972, October~2003,
January~2005; parameters in Supplementary Table~S3). BFO absorbed dose is
computed via CSDA transport through wall plus $10\gcm$ tissue and compared
to the NASA 30-day BFO limit of $250$~mGy \citep{cucinotta2010}. Mission SEP
encounter probability follows a Poisson process at solar-cycle-dependent
rates \citep{feynman1993}.

\subsection{Uncertainty Quantification}
\label{sec:uq}

Nine parameters are varied simultaneously via LHS: five physics
(solar modulation, LIS normalization, nuclear
cross-sections, neutron yield, HZE composition) and four biology (quality
factor scale, DDREF, ERR scale, EAR scale), each with a lognormal, uniform,
or clipped-normal prior justified by the relevant literature spread
(full distributions and justifications in Supplementary Table~S4). For each
LHS sample, a full mission dose integration is performed; publication
results use $N=500$ samples. Parameter sensitivity uses normalized
Spearman-$\rho^2$ decomposition (an approximation to first-order Sobol
indices).
\label{sec:sobol}

A separate value-of-information sweep (\S\ref{sec:vov}) re-runs the $N=500$
ensemble once per parameter, halving that parameter's prior spread while
holding all others at baseline, to quantify how much narrowing any one
source of uncertainty would actually reduce REID uncertainty
(\code{narrow} argument to \code{gcrisk.uncertainty.lhs\_samples}; Supplementary
\S{}S1.9).
\label{sec:vovmethod}

\subsection{Dose-Averaged LET and Proton RBE Self-Consistency Check}
\label{sec:letmethod}
\label{sec:rbemethods}

Per-species and field dose-averaged LET ($\overline{L}_D$) are computed to
trace $\Qeff$ to specific ion contributions (\S\ref{sec:let}; full
definitions in Supplementary \S{}S1.9). The same $\overline{L}_D$
infrastructure also implements two proton RBE models, Wedenberg et~al.\
(2013) \citep{wedenberg2013} and McNamara et~al.\ (2015)
\citep{mcnamara2015}, matched to OpenTOPAS-RBE scorer conventions and
checked only as an implementation self-consistency test
(Table~\ref{tab:validation}). This
layer is the foundation for a separate, planned proton-therapy crossover
study.
